\section{目标后端：从 LLVM IR 到 RV64GC 机器程序}
\label{sec:backend}

中端的输出仍拥有无限个虚拟 SSA 值、抽象操作和结构化的类型；真实处理器只有有限寄存器、特定指令编码、调用约定和流水线资源。后端的任务不是把每条 IR 指令替换成一条汇编，而是在目标语义与成本约束下共同解决操作合法化、指令选择、调度、寄存器分配、栈帧实现和代码发射。LLVM 官方代码生成器文档给出的主干是：指令选择与调度形成目标机器指令，在物理寄存器分配前后进行机器级优化，随后插入序言/尾声并经 MC 层发射汇编或对象文件\cite{llvm-codegen}。具体 pass 顺序会随目标、版本和优化级别改变，不能把这条概念主干误写为固定流水线。

本节以仓库根目录下实际捕获命令的等价写法
\begin{lstlisting}[basicstyle=\ttfamily\scriptsize]
clang --target=riscv64-unknown-linux-gnu -march=rv64gc -mabi=lp64d \
  -std=c17 -Ipreflight/include -O2 -fno-inline -S \
  preflight/src/clipped_dot.c \
  -o .temp/preflight-rv64/artifacts/clipped_dot.O2.s
\end{lstlisting}
及 Clang 18.1.3 生成的汇编为观察对象。\texttt{rv64gc} 表示 64 位基础整数 ISA 加通用扩展集合，ABI 为 LP64D。工件中没有保存指令选择器标识、逐 pass MIR 或调度日志；因此对最终汇编的逐条分析是\observed{}，而“SelectionDAG 还是 GlobalISel 产生某条指令”不会被冒充为观察事实。

\subsection{后端不是一张查找表}

LLVM IR 允许目标不直接支持的整数位宽、向量形状、原子操作、内在函数和调用形式。合法化（legalization）先把操作和类型拆分、扩展、收窄或改写为目标能够承担的形式，必要时还会生成运行库调用。随后，指令选择（instruction selection）在等价模式之间寻找目标指令，并利用立即数、寻址模式和融合机会。概念数据流可写成：

\begin{center}
\small\ttfamily
LLVM IR
$\rightarrow$ lowering/legalization
$\rightarrow$ selected MachineInstr
$\rightarrow$ scheduling
$\rightarrow$ register allocation
$\rightarrow$ frame lowering
$\rightarrow$ MC
\end{center}

这个箭头图刻意省略了多次组合、前后调度和机器级清理；它说明数据依赖，而不承诺单一线性 pass 表。

LLVM 目前有两条主要选择框架。SelectionDAG 把一个基本块附近的操作与依赖构造成有向无环图（Directed Acyclic Graph, DAG），经历 DAG 组合、类型/操作合法化、目标模式匹配和调度，最后形成 \texttt{MachineInstr}\cite{llvm-codegen}。GlobalISel 则明确分为：
\begin{center}
\small\ttfamily
LLVM IR $\rightarrow$ IRTranslator $\rightarrow$ generic MIR
$\rightarrow$ Legalizer $\rightarrow$ RegisterBankSelect
$\rightarrow$ InstructionSelect $\rightarrow$ target MIR
\end{center}
其中通用机器 IR（generic Machine IR, gMIR）先保留 \texttt{G\_ADD} 等通用操作，合法化后分配寄存器银行，再消除通用操作\cite{llvm-globalisel-pipeline,llvm-globalisel}。GlobalISel 文档仍明确提示支持程度与回退路径具有目标差异。两条路线可以生成相同的最终汇编；仅凭 \texttt{.s} 无法反推实际选择器。若要回答这一工程问题，应保存 Clang/\texttt{llc} 的明确选择器选项、调试日志或 MIR 检查点，而不是依据目标名称猜测。

机器 IR（Machine IR, MIR）位于 LLVM IR 与最终汇编之间。它含机器基本块、目标或通用机器指令、虚拟/物理寄存器、live-in、栈帧对象和 pass 属性；可序列化为供 \texttt{llc -run-pass} 测试使用的 YAML 文本\cite{llvm-mir}。MIR 是极有价值的后端测试接口，却不是稳定的用户 ABI；本项目未保存 MIR，因此本节用它解释“阶段之间传递什么”，不展示虚构的案例转储。

\subsection{实例映射一：循环被改写为指针归纳}

生成的 \texttt{clipped\_dot} 入口遵循 RISC-V ELF psABI：\texttt{a0}/\texttt{a1} 是两个数组地址，\texttt{a2} 为 \texttt{n}，\texttt{a3}/\texttt{a4} 是上下界，整数返回值置于 \texttt{a0}\cite{llvm-riscv-psabi}。表~\ref{tab:rv64-map} 沿实际汇编追踪关键 IR 语义。

\begin{table}[t]
  \centering
  \caption{优化 IR 语义到实际 RV64GC 汇编的映射}
  \label{tab:rv64-map}
  \scriptsize
  \begin{tabular}{p{0.28\textwidth}p{0.30\textwidth}p{0.32\textwidth}}
    \toprule
    IR/源语义 & 工件中的 RV64 片段 & 后端决定 \\
    \midrule
    \(n\le0\) 时零次迭代 & \texttt{li a5,0; blez a2,.LBB0\_7} & 累加器预置 0，直接复用出口 \\
    循环上界与索引 & \texttt{slli a2,a2,2; add a6,a1,a2} & 把 \texttt{i<n} 改为 \texttt{b\_ptr==end}，消除显式 \texttt{i} \\
    两个 \texttt{i32} load & \texttt{lw a7,0(a0); lw a2,0(a1)} & 4-byte load；在 RV64 上符号扩展为寄存器宽度 \\
    \texttt{mul i32} & \texttt{mulw a2,a2,a7} & W 后缀保留低 32 bit 并符号扩展结果 \\
    下/上界截断 & \texttt{blt}、\texttt{mv} 与汇合块 & 选择分支实现，而非要求目标拥有 min/max 指令 \\
    累加与下一元素 & \texttt{addw a5,t0,a5}; 两个 \texttt{addi ...,4} & 32-bit 求和；指针归纳替代每轮 \(i\times4\) \\
    循环退出与返回 & \texttt{beq a1,a6,...; mv a0,a5; ret} & 端指针比较；把结果放入 ABI 返回寄存器 \\
    \bottomrule
  \end{tabular}
\end{table}

核心片段为：
\begin{lstlisting}[language=RISCV,basicstyle=\ttfamily\scriptsize]
li      a5, 0
blez    a2, .LBB0_7
slli    a2, a2, 2
add     a6, a1, a2          # end = b + 4*n
...
lw      a7, 0(a0)
lw      a2, 0(a1)
mulw    a2, a2, a7
...
addw    a5, t0, a5
addi    a1, a1, 4
addi    a0, a0, 4
beq     a1, a6, .LBB0_7
\end{lstlisting}

这里发生了一个值得借鉴的“换角度”：源代码用整数下标 \texttt{i}，手写 IR 用 \texttt{sext+GEP}，机器代码却保留两个不断前进的指针和一个末端指针。这样每轮不再计算 \(i\times4\)，循环退出也只需一次地址比较。它是归纳变量简化、地址模式选择和目标成本共同作用后的形状，不能机械归因给单条 \texttt{getelementptr}。

RV64 寄存器为 64 位，但 C/SysY \texttt{int} 与 IR \texttt{i32} 仍是 32 位。\texttt{lw} 将 32 位值符号扩展到 XLEN；\texttt{mulw} 和 \texttt{addw} 只计算低 32 位，再把 bit 31 符号扩展到 64 位寄存器。后端因此不需要给每个 \texttt{i32} SSA 值分配“物理 32 位寄存器”；同一物理寄存器承载的值语义由生成该值的指令及后续使用决定。这也是合法化和选择必须与目标寄存器类协同的原因。

\subsection{实例映射二：调用、常量地址与栈帧}

\texttt{clipped\_dot} 是叶函数（leaf function）：它不调用其他函数，所有活跃值能放入调用者保存的整数寄存器，因此实际工件没有调整 \texttt{sp}，也没有保存 \texttt{ra}。\texttt{ret} 是对 RISC-V 返回序列的汇编表示。相反，\texttt{main} 调用四个函数，生成：
\begin{lstlisting}[language=RISCV,basicstyle=\ttfamily\scriptsize]
addi    sp, sp, -16
sd      ra, 8(sp)
call    getint
...
call    clipped_dot
call    putint
li      a0, 10
call    putch
...
ld      ra, 8(sp)
addi    sp, sp, 16
ret
\end{lstlisting}
16-byte 栈调整满足调用边界的栈对齐，保存 \texttt{ra} 是因为任意 \texttt{call} 会覆盖返回地址。psABI 把 \texttt{ra} 归为调用者保存寄存器；此处恰是 \texttt{main} 作为调用者，为跨越其内部调用而保存自己的返回地址，并非 \texttt{s0--s11} 那类“被调用者必须替来访者保存”的义务。该 \texttt{sd}/\texttt{ld ra} 也\emph{不是}寄存器分配失败造成的普通 spill。把所有栈读写都叫“溢出”会掩盖两个不同机制：一个保护跨调用仍要使用的架构状态，另一个才是虚拟值竞争物理寄存器失败后产生的临时槽。

\texttt{-O2} 已将两个初始化数组保留在只读数据区，\texttt{main} 用 PC 相对序列取址：
\begin{lstlisting}[language=RISCV,basicstyle=\ttfamily\scriptsize]
.Lpcrel_hi0:
auipc   a0, %pcrel_hi(.L__const.main.a)
addi    a0, a0, %pcrel_lo(.Lpcrel_hi0)
\end{lstlisting}
\texttt{auipc} 与随后 \texttt{addi} 共同构成地址；\texttt{\%pcrel\_hi/lo} 仍需汇编器和链接器用重定位确定，汇编文本并非最终数值编码。传参随后复用 \texttt{a0}/\texttt{a1}，\texttt{a3=-8}、\texttt{a4=12}。调用返回后，结果已经位于 \texttt{a0}，所以紧接的 \texttt{call putint} 无需搬移。这是指令选择与调用约定一起消除拷贝的例子。

输入长度的 \([0,8]\) 截断在汇编中先处理上界，再用布尔值构造 0 或原值：
\begin{lstlisting}[language=RISCV,basicstyle=\ttfamily\scriptsize]
li      a1, 8
blt     a0, a1, .LBB1_2
li      a0, 8
.LBB1_2:
sgtz    a1, a0
neg     a2, a1
and     a2, a2, a0
\end{lstlisting}
若 \texttt{a0>0}，\texttt{sgtz} 产生 1，\texttt{neg} 得到全 1 掩码，于是 \texttt{a2=a0}；否则掩码为 0。源 CFG、优化 IR 的 \texttt{smin/smax} 与这段分支加掩码代码语法完全不同，但在整数域上承担同一值函数。后端正确性的评价对象应是语义关系，而非逐行相似度。

\subsection{调度：依赖、延迟和寄存器压力之间的折衷}

指令调度（instruction scheduling）要在线性顺序中满足真依赖、反依赖以及目标 hazard，同时考虑乘法/load 延迟、发射资源和代码尺寸。SelectionDAG 路线可在形成 \texttt{MachineInstr} 前调度 DAG；某些目标还有寄存器分配前后调度或后期 hazard 修复。因此“选择---分配---调度”不是所有后端唯一的时间顺序\cite{llvm-codegen}。

案例循环中两次独立 load 可以在乘法之前发出，指针更新则必须在本轮 load 取走旧地址以后语义上生效。把更多独立操作提前可能隐藏延迟，却会延长值的活跃区间（live range），增加寄存器压力；反过来，为减少压力而紧缩活跃区间可能损失指令级并行。最终汇编显示了一个合法顺序，但没有微架构模型和调度日志，不能仅凭肉眼断言它是“最优调度”，更不能确定是哪一个调度器做出决定。

\subsection{寄存器分配与 spill：有限资源如何承载无限 SSA 名称}

寄存器分配（register allocation）把虚拟寄存器的活跃区间映射到有限且存在别名关系的物理寄存器，并服从预着色参数、调用破坏集合和寄存器类限制。经典图着色模型把同时活跃的值连为干涉图，相邻结点不得取同一颜色\cite{llvm-chaitin1981}；LLVM 的优化分配器则还使用 live interval、拆分、驱逐和 spill 成本等启发式，不能简化为“直接跑图着色”\cite{llvm-codegen}。

若可用颜色不足，分配器可把值溢出（spill）到栈槽，在使用前 reload；必要时还会重物化（rematerialization），即重新计算廉价常量而非读取栈。案例内核压力较低，\observed{}的函数体没有 \texttt{sp} 访问，故可说“最终内核没有 spill”；但这不证明分配过程从未尝试过其他方案。寄存器复用很积极，例如输入 \texttt{n} 原在 \texttt{a2}，计算完 \texttt{end} 后 \texttt{a2} 被复用为数组元素和乘积，因为原 \texttt{n} 已死亡。SSA 名称无限，并不意味着每个名称拥有永久物理寄存器。

对含调用的 \texttt{main}，调用者保存寄存器（caller-saved register）上的临时值若需跨调用存活，必须在调用前保存或被分配到被调用者保存寄存器；被调用者保存寄存器（callee-saved register）一旦使用，则由当前函数在序言/尾声中恢复。实际 \texttt{main} 足够简单，仅保存 \texttt{ra}，数组已在常量区，因而帧大小只有 16 byte，而不是手写汇编为两个栈数组准备的 80 byte。优化改变了对象存储策略，帧降低（frame lowering）据此实现具体偏移；栈帧大小不是仅由源局部变量数决定的。

\subsection{序言/尾声、伪指令与晚期机器优化}

帧降低发生在 spill 槽、被保存寄存器和对齐需求已知之后，它把抽象 frame index 解析为 \texttt{sp}/frame-pointer 相对地址，插入栈调整、寄存器保存恢复以及调试/展开所需的 CFI 指令。案例汇编中的 \texttt{.cfi\_def\_cfa\_offset} 与 \texttt{.cfi\_offset} 不执行算术，而为异常展开和调试器描述帧状态。

汇编还含 \texttt{li}、\texttt{mv}、\texttt{ret}、\texttt{call} 等伪指令（pseudo-instruction）。它们在后续汇编阶段根据立即数范围、目标扩展和重定位展开为一个或多个真实编码。例如 \texttt{mv} 常可编码为 \texttt{addi rd,rs,0}，\texttt{ret} 通常对应通过 \texttt{ra} 的间接跳转；\texttt{call} 可能形成带重定位的一组指令，并在链接时进一步松弛。故“汇编一行等于一条机器指令”也不是稳定假设。

机器级 peephole、分支折叠、拷贝合并和死指令删除会继续清理选择与分配留下的冗余。例如 \texttt{phi} 在进入非 SSA 机器形式前通常被实现为前驱边上的并行拷贝，再通过关键边拆分和 coalescing 尽量消去；最终汇编当然不会出现 \texttt{phi} 操作码。这一消失并非源信息丢失，而是把“按来边选值”的抽象约束落实为控制流边上的寄存器状态。

\subsection{MC 层：同一个 MachineInstr，可走向文本或对象字节}

寄存器分配和晚期机器 pass 完成后，LLVM 把 \texttt{MachineInstr} 降为 MC 层的 \texttt{MCInst} 与操作数。\texttt{MCStreamer} 接收标签、节、数据、指令和指示；\texttt{MCAsmStreamer} 输出 \texttt{.s} 文本，而 \texttt{MCObjectStreamer} 配合 \texttt{MCCodeEmitter} 直接编码对象文件并产生 fixup/重定位\cite{llvm-codegen}：
\begin{center}
\small\ttfamily
MachineInstr $\rightarrow$ MCInst
$\begin{cases}
\rightarrow \text{MCAsmStreamer} \rightarrow .s \rightarrow \text{assembler} \rightarrow .o\\
\rightarrow \text{MCObjectStreamer+MCCodeEmitter} \rightarrow .o
\end{cases}$
\end{center}

因此 \texttt{clang -S} 保存汇编只是一个可观察选择；\texttt{clang -c} 使用集成汇编器时可以不在磁盘上产生任何 \texttt{.s}。对象文件也不只是指令字节：它还必须携带节、符号、字符串表、对齐、重定位以及可选调试/展开信息。上面的 \texttt{\%pcrel\_hi/lo} 和 \texttt{call getint} 在目标地址未知时正需要这种结构。汇编器、链接器和 ELF 的后续责任将在下一阶段继续展开。

\subsection{我们已经知道什么，还不能知道什么}

表~\ref{tab:backend-evidence} 汇总本节的认识边界。其目的不是保守措辞，而是让每项结论都有可复查证据类型。

\begin{table}[t]
  \centering
  \caption{后端结论的证据边界}
  \label{tab:backend-evidence}
  \small
  \begin{tabular}{p{0.42\textwidth}p{0.42\textwidth}}
    \toprule
    可由保存工件直接确认 & 需要新增日志/实验才能确认 \\
    \midrule
    目标为 \texttt{riscv64-linux-gnu}、\texttt{rv64gc/lp64d}；最终为标量循环 & 实际采用 SelectionDAG、GlobalISel 还是回退路径 \\
    参数/返回寄存器、叶函数无栈帧、\texttt{main} 的 16-byte 帧 & 每条机器指令由哪个具体 pass/匹配规则产生 \\
    内核最终无 spill，\texttt{main} 保存 \texttt{ra} & 调度器考虑的精确微架构延迟与候选顺序 \\
    \texttt{i32} 使用 \texttt{lw/mulw/addw}；数组遍历使用指针归纳 & 中间 MIR、live interval、分配器驱逐/拆分历史 \\
    汇编含 PC 相对取址、符号调用和 CFI & 目标文件中的最终编码与重定位，须检查 \texttt{readelf/objdump} \\
    \bottomrule
  \end{tabular}
\end{table}

最终汇编可由独立汇编器接受并链接运行，是重要的端到端证据；但它不能证明后端每个中间决策，也不能穷尽输入等价性。后端验证的可操作路线是：用 MIR 单元测试隔离机器 pass，用 \texttt{llvm-mca} 或硬件计数器检验调度/吞吐假设，用 \texttt{readelf}/\texttt{objdump} 核对编码与重定位，并以差分测试连接源、IR、手写汇编和最终可执行文件。这样，性能主张、结构主张和语义主张就不会被一张看似漂亮的汇编清单混在一起。
